\documentclass[conference]{IEEEtran}
\usepackage[utf8]{inputenc}
\usepackage[T1]{fontenc}
\usepackage{cite}
\usepackage{amsmath,amssymb}
\usepackage{graphicx}
\usepackage{booktabs}
\usepackage{url}

\title{Can Deterministic Network Verification Reduce Undetected Faults in LLM‑Generated Configurations?}

\author{
\IEEEauthorblockN{Autonomous AI Researcher}
\IEEEauthorblockA{CNSM 2026 Agentic AI Researcher Track}
}

\begin{document}

\maketitle

\begin{abstract}
We preregistered and executed a paired confirmatory study (study id c1-gnr-vv-2026-0001) to measure whether inserting a deterministic network verifier between an LLM intent\ensuremath{\rightarrow}configuration generator and an emulated execution reduces the proportion of configurations that would execute with operational faults. Using shared\_initial\_candidate semantics (one identical LLM candidate per intent evaluated both without and with verification), we evaluated 72 preregistered paired intents with gpt-5-mini and a canonical deterministic verifier (deterministic\_netops\_validator\_v5). Execution completed with 144 episodes (72 pairs) and 72 model calls. All baseline executions produced no emulator-observed faults (n\_11 = 72, n\_10 = n\_01 = n\_00 = 0). The paired difference in undetected-fault proportion is 0.0 (paired bootstrap 95\% CI [0.0, 0.0], exact McNemar two-sided p = 1.0; bootstrap seed = 1729, resamples = 10000). Under the preregistered shared\_initial\_candidate contract this degenerate confirmatory outcome is expected when identical generated candidates produce no baseline execution faults. We report the preregistration rationale, deterministic execution accounting, representative validator diagnostics, exact inferential outputs, and artifact-grounded recommendations for follow-on confirmatory designs able to detect verifier utility under realistic baseline-error regimes.
\end{abstract}

\section{Introduction}
Generative large language models (LLMs) are increasingly proposed to translate high-level operator intents into device-level network configurations. This intent\ensuremath{\rightarrow}configuration automation can speed change pipelines but risks silently violating reachability, access-control, ordering, or workflow invariants and thereby causing outages if generated text is executed without additional checks. A standard operational mitigation is a generate--verify--repair (GVR) architecture that combines stochastic LLM generation with deterministic verification and, if needed, repair or human review. However, despite many architectural proposals and prototypes, systematic preregistered experiments that measure a verifier's detection performance against executed outcomes under tightly controlled paired designs are scarce.

This paper answers a narrowly scoped, preregistered research question: does inserting a deterministic network verifier between an LLM generator and emulator execution materially reduce the proportion of configurations that execute with operational faults when the identical generated candidate is evaluated both without and with verification? That estimand isolates verifier detection from improvements that could arise from re-sampling, prompt-engineering, or repair-enabled iterations. We executed study c1-gnr-vv-2026-0001 using gpt-5-mini and deterministic\_netops\_validator\_v5 across 72 preregistered paired intents under shared\_initial\_candidate semantics. The run completed with no technical failures, produced a degenerate confirmatory outcome (no baseline emulator faults), and yields a paired difference of 0.0 (95\% CI [0.0,0.0], exact McNemar p = 1.0). Rather than treating this as a null failure, we analyze why it occurred given the preregistered contract and archived artifacts, and we provide concrete, artifact-grounded guidance for designs that can detect verifier utility when baseline error prevalence is non-negligible.

\section{Related Work}
We position this study against three closely related strands: (1) LLM‑based intent\ensuremath{\rightarrow}configuration automation and emulated evaluations; (2) multi-agent and tool‑augmented NetOps frameworks that propose verifier integration; and (3) generate--verify--repair and constrained‑decoding methods from software/code domains that inform possible NetOps adaptations.

LLM intent\ensuremath{\rightarrow}configuration translation. Recent empirical work demonstrates that instruction‑tuned LLMs can map operator intents to device‑level configuration text within constrained vendor dialects and curated testbeds. For example, Nguyen et al. present an intent\ensuremath{\rightarrow}configuration pipeline and evaluate LLM performance on structured configuration tasks, documenting promising task‑level accuracy but highlighting the gap between textual correctness and safe execution in operational environments [1]. Those studies use emulator or synthesized workloads to assess functional correctness of generated artifacts but emphasize that textual matching metrics do not substitute for execution‑level invariants such as reachability or ordering constraints. Our study adopts the same operational framing (emulator execution as a conservative proxy for operational fault) but differs methodologically by preregistering a paired detection estimand under shared\_initial\_candidate semantics to isolate verifier detection capability on identical generated candidates.

Multi-agent and verifier‑integrating NetOps systems. Several prototypes and architectural reports advocate tool‑augmented or multi‑agent LLM systems that incorporate verification steps to reduce regressions and to enable safe automation at scale. Notably, recent multi‑agent intent‑driven frameworks demonstrate how separate agent roles (planner, config generator, validator) can be orchestrated to produce and check configurations before execution; these systems show how verifier signals can be folded into human‑in‑the‑loop or automated repair workflows but typically report system behavior at the demonstration or prototype level rather than in preregistered paired detection trials [2]. Our experiment operationalizes a single, canonical verifier in a strict generate--verify evaluation contract to measure whether the verifier can, by itself and without changing the generated candidate, prevent executed faults that would otherwise occur.

Generate--Verify--Repair and constrained decoding provenance. The generate--verify--repair (GVR) pattern has a substantial methodological pedigree in program repair and structured‑output generation. Architectural proposals and empirical studies in code generation and self‑repair show that deterministic checkers (unit tests, static analyzers) combined with stochastic generation can enable iterative repair loops that converge on functionally correct outputs [3]. Complementary work on constrained decoding and integer‑programming‑based decoding demonstrates that enforcing structural constraints during generation reduces syntactic or semantic violations in structured domains [4]. These results motivate two dimensions of adaptation for NetOps: (a) constraining or biasing generation toward verifier‑satisfiable forms (constrained decoding), and (b) allowing verifier‑triggered repair iterations that change the executed candidate distribution. Importantly, those adaptations change the causal estimand: they do not measure verifier detection on an identical candidate but rather the end‑to‑end benefit of verifier‑enabled generation and repair. The preregistered shared\_initial\_candidate semantics used here intentionally hold the generated candidate fixed across conditions so that any paired difference can be attributed only to verifier detection (not to sampling, constrained decoding, or repair).

Deterministic network verification and policy checking. Deterministic analysis tools developed in the networking literature (header‑space analysis and related formal checkers) provide precise, semantics‑aware invariants for reachability, ACL, and forwarding behavior; these tools are natural verifier candidates for LLM‑generated configs because they can conservatively flag violations that imply operational harm [5]. Prior work on real‑time policy checking and emulated validation shows how formal analyses can be integrated into testbeds to prevent regressions before deployment. Our canonical verifier (deterministic\_netops\_validator\_v5) was used in the same role: apply deterministic, rule‑based checks that compare the generated sequence against manifested workflow constraints (sequence, allowed fields, and protected interfaces) encoded in each task manifest. Where prior literature argues for verifier integration at the architecture level, our contribution is a preregistered paired measurement that quantifies verifier detection power when the generated candidate is held identical across conditions.

Synthesis and gap. Collectively the literature indicates (a) LLMs can produce plausible device commands under constrained settings, (b) multi‑agent and tool‑augmented systems often propose a verifier role, and (c) GVR and constrained decoding strategies can materially change generation outcomes by construction. What is missing---and what this study targets---is a preregistered, paired, artifact‑grounded measurement of verifier detection performance on identical generated candidates: that specific estimand isolates the marginal detection capacity of a verifier decoupled from generation modifications. The degenerate result we observe (zero baseline emulator faults) is thus directly informative: it reveals a design regime---highly constrained synthetic intents and generator behavior---where detection cannot be measured because baseline faults are absent. We use archived traces and validator diagnostics to show how task specification and verifier rule alignment created that regime, and we suggest preregistered alternatives (independent sample arms, repair flows, adversarial task banks) that the literature suggests would expose measurable verifier utility when baseline error prevalence is non-negligible.

\section{Methodology}
Preregistration and primary estimand. We preregistered and froze study c1-gnr-vv-2026-0001 before observing outcomes. The primary estimand is the paired reduction in undetected operational-fault proportion (paired\_success\_rate\_difference\_guarded\_minus\_baseline). The confirmatory hypothesis (H1) targeted an absolute paired reduction of 0.20 (two-sided alpha = 0.05, power >= 0.80); a sample-size heuristic yielded \ensuremath{\approx}63 pairs and we preregistered N = 72 pairs (24 per topology-difficulty stratum). The preregistration and analysis plans, including sensitivity analyses and exclusion rules, are part of the archived preregistration record.

Shared\_initial\_candidate semantics and experimental contract. The execution\_contract enforced shared\_initial\_candidate semantics: for each intent we performed exactly one initial LLM generation and evaluated that identical candidate in two paired conditions. Baseline condition: execute the shared candidate in an isolated emulator and record whether emulator-observed operational faults occur. Guarded condition: run the canonical deterministic verifier on the same candidate; if the verifier flags violations the guarded outcome is recorded as prevented (no executed fault); if the verifier does not flag, execute the identical candidate in the emulator and record the result. By construction, any observed paired difference can only arise from the verifier preventing execution of a candidate that would otherwise have failed at emulation; this isolates detector utility from generator-side resampling or repair effects.

Model, adapter, and verifier configuration (artifact-grounded). The declared model is gpt-5-mini (execution/model\_configuration.json), requested via the hosted\_netops\_gvr\_v1 adapter. Key recorded model parameters: declared\_model\_version = "gpt-5-mini", provider = "openai\_responses", max\_output\_tokens = 2000, maximum\_attempts\_per\_call = 1, and temperature = null. The recorded temperature = null indicates that provider-default runtime sampling semantics were used; we treat this as an archived sampling-parameter uncertainty rather than an explicit numeric temperature. The canonical verifier is deterministic\_netops\_validator\_v5 (validator\_version = "5.0") configured to run the full\_intent\_constraint\_validation rule set. The verifier emits structured validator\_traces per episode (fields include parsed\_command\_count, required\_command\_count, normalized\_configuration, validation\_before and validation\_after final\_state snapshots, violation\_codes, violation\_count, and difficulty metadata). The verifier's validity verdict and trace for each episode were archived under execution/scoring/ and formed the mechanistic basis for guarded decisions.

Task-bank construction, contamination controls, and stratification. Confirmatory tasks were generated deterministically by deterministic\_netops\_task\_generator\_v5. Each task\_manifest entry encodes: a natural-language intent, a machine-structured required\_sequence (ordered field changes that implement the intent), allowed\_touched\_fields, preserved-state policies, initial and expected\_final\_state snapshots, and a difficulty.level tag (medium/hard/very\_hard). The preregistration required deterministic exact-match contamination checks against a public-corpus fingerprint (analysis/contamination\_summary.csv); exact-match flagged items are excluded from confirmatory analysis (none were excluded in this run). The confirmatory sample followed the preregistered stratification (24 pairs per difficulty stratum) to allow descriptive subgroup inspection.

Execution protocol, retries, and deterministic accounting. The missingness\_plan permitted up to two retries (three attempts total) for transient hosted-API errors on generation calls; all retries and provider events were logged. The run started at 2026-08-31T06:56:23.065707Z and completed at 2026-08-31T07:06:28.665518Z (execution manifest). The execution manifest records planned\_episode\_count = 144, completed\_episode\_count = 144, failed\_episode\_count = 0, and model\_calls\_used = 72 (one shared generation per pair). Provider-call audit (deterministic\_reconciliation.json) reports hosted\_model\_call\_event\_count = 72, completed\_hosted\_model\_call\_event\_count = 72, failed\_hosted\_model\_call\_event\_count = 0, cache\_miss\_count = 72, and stage\_counts = \{"shared\_initial\_generation": 72\}. These deterministic accounting artifacts were used to reconcile paired outcomes and to confirm that the shared\_initial generation stage was the only cross-condition generation activity.

Scoring, validation rules, and per-episode traces. For each episode the verifier produced a binary validity verdict (valid / invalid) and an accompanying validator\_trace JSON that documents parsed and required command counts, normalized\_configuration, validation\_before/after states, and any violation\_codes. The scoring rule deployed was full\_intent\_constraint\_validation (validator enforces required\_sequence order, allowed\_touched\_fields, and preserve\_unspecified\_state constraints encoded in the task manifest). Per-episode scoring artifacts (execution/scoring/task-*.json) therefore permit audit of why a candidate was accepted or flagged. In the recorded run no verifier flagging occurred for confirmatory pairs (violation\_count = 0 across representative episodes), and no automated repair was applied (repair\_applied = false in validator\_trace fields).

Inferential and sensitivity analysis procedures (deterministic scripts). The primary analysis used the registered paired\_binary\_analysis\_v1 executor to construct the 2\ensuremath{\times}2 contingency table (guarded \ensuremath{\times} baseline) using binary indicators where an undetected executed fault is 1 and success/no-fault is 0. The prespecified exact inferential test is an exact McNemar two-sided test when discordant pairs exist. We also computed paired bootstrap-percentile 95\% confidence intervals for the paired difference using 10,000 resamples (bootstrap\_seed = 1729), recorded in analysis/analysis\_log.jsonl. Pre-specified subgroup confirmatory comparisons (topology complexity and policy type) were registered with Holm correction to control familywise error; these controls are implemented in the deterministic analysis pipeline but become non-informative in the absence of discordant pairs. Pre-registered sensitivity analyses executed by the same scripts include (a) treating excluded confirmatory pairs as failures (complete\_pair\_primary\_with\_failure\_as\_zero\_sensitivity) and (b) bootstrap diagnostics for detector detection and false-positive rates. Because the run produced no exclusions and no discordant outcomes these sensitivity analyses reproduce the degenerate numerical conclusion.

Design-level notes and construct tradeoffs (artifact-linked). The shared\_initial\_candidate semantics intentionally isolates verifier detection from generator-side sampling; this preserves a clean causal interpretation of verifier-prevention utility but prevents measurements that rely on verifier-triggered resampling or repair. The archived execution/model\_configuration.json recording temperature = null documents that provider-default sampling behavior was allowed; had we required explicit sampling parameters (non-null temperatures) we could have measured sampling-sensitivity as a separate factor. The verifier's full\_intent\_constraint\_validation rule set enforces task-manifest constraints: because the task generator encoded explicit required\_sequence and allowed\_touched\_fields many generated candidates closely matched those constraints, which the verifier therefore did not flag. This alignment is measurable in validator\_trace fields (parsed\_command\_count equals required\_command\_count and violation\_count = 0 for representative episodes) and explains the ceiling result when baseline emulator faults are absent. The preregistered sensitivity and alternative contracts (independent-generation arms, repair-enabled flows, adversarial task seeding, and multi-model comparisons) are described in the preregistration and can be executed using the same deterministic execution and analysis infrastructure archived with this run.

\section{Results}
Execution and manifest summary. The autonomous pipeline completed without technical failures. Key manifest figures (execution/ and analysis/ manifests): planned\_episode\_count = 144; completed\_episode\_count = 144; failed\_episode\_count = 0; model\_calls\_used = 72; artifact\_count = 510. analysis/condition\_summary.csv reports baseline successes = 72, baseline failures = 0 (success\_rate = 1.0) and guarded successes = 72, guarded failures = 0 (success\_rate = 1.0).

Paired contingency and exact inference. Aggregating across the 72 confirmatory pairs produced contingency counts n\_11 = 72, n\_10 = 0, n\_01 = 0, n\_00 = 0 (guarded \ensuremath{\times} baseline). The paired estimate paired\_success\_rate\_difference\_guarded\_minus\_baseline = 0.0. Exact McNemar two-sided p = 1.0. The paired bootstrap percentile 95\% CI (seed = 1729; 10,000 resamples) is [0.0, 0.0], reflecting the degenerate distribution produced by the absence of discordant pairs. analysis/paired\_contingency\_table.csv and analysis/condition\_summary.csv contain these tabulations and are archived.

Representative validator diagnostics and per-episode exemplars. Representative per-episode scoring artifacts are archived (execution/scoring/task-000001-*.json ... task-000006-*.json). Representative summaries (extracted from archived JSONs):
- pair-000001 (task-000001): difficulty.level = "medium", parsed\_command\_count = 4, required\_command\_count = 4, normalized\_configuration matched the shared candidate, validator\_version = "5.0", violation\_count = 0.
- pair-000002 (task-000002): difficulty.level = "hard", parsed\_command\_count = 2, required\_command\_count = 2, violation\_count = 0.
- pair-000003 (task-000003): difficulty.level = "hard", parsed\_command\_count = 6, required\_command\_count = 6, violation\_count = 0.
These traces show the verifier executed on each shared candidate and produced no violation flags across the sampled representative tasks. Execution/scoring/*.json files contain full validator\_traces for reproducibility.

Contamination, missingness, and sensitivity diagnostics. analysis/contamination\_summary.csv reports zero exact-match contaminations for confirmatory pairs. analysis/missingness\_summary.csv records complete\_pairs = 72 and zero failed or unscorable episodes. Pre-specified sensitivity analyses (treating excluded pairs as failures; bootstrap diagnostics) were executed and archived; because no pairs were excluded and no discordant outcomes occurred these sensitivity analyses reproduce the degenerate numeric conclusion (paired difference = 0.0). analysis/analysis\_log.jsonl records bootstrap\_seed = 1729 and bootstrap\_resamples = 10000.

\section{Discussion}
Confirmatory interpretation and boundary conditions. The preregistered paired design with shared\_initial\_candidate semantics validly measures verifier detection on identical candidates; because the identical generated candidates produced zero emulator faults the verified paired outcome is a ceiling/null result: the verifier could not reduce executed faults that did not occur. The numeric outputs (n\_11 = 72; paired difference = 0.0; McNemar p = 1.0; bootstrap CI [0.0,0.0]) are direct consequences of the preregistered contract and observed data and therefore correctly reported as a degenerate confirmatory result.

Artifact-grounded causes of the ceiling outcome. Archived artifacts allow us to identify measurable factors that explain the ceiling outcome: (1) the deterministic task generator produced constrained, high-clarity intents with explicit required\_sequence fields and allowed\_touched\_fields encoded in each task manifest, increasing the likelihood of unambiguous correct candidates; (2) gpt-5-mini (declared\_model\_version in execution/model\_configuration.json) generated sequences that matched those structured constraints across the synthetic corpus; and (3) the verifier rule set (full\_intent\_constraint\_validation in deterministic\_netops\_validator\_v5) aligns closely with the required\_sequence constraints present in the task manifests, making verification and generation objectives strongly congruent. These archived, measurable factors together account for the absence of baseline emulator faults.

Design implications: when verification benefit is measurable. To empirically measure practical verifier benefits, confirmatory contracts must present non-negligible baseline error prevalence or allow verifier-mediated candidate changes. Artifact-grounded, preregisterable alternatives include: (A) independent-generation arms---allow separate initial generations per condition so verifier-triggered repairs or alternative sampling may change executed candidate distributions; (B) repair-enabled flows---permit a deterministic repair call when the verifier flags violations and execute repaired candidates to measure end-to-end improvement; (C) adversarial/fault-injection task banks---seed tasks with ambiguous intents, ordering subtleties, or vendor-edge cases to raise baseline error prevalence; (D) multi-model and sampling sweeps---include multiple model versions and explicit sampling parameter settings (non-null temperatures) to quantify model- and sampling-sensitivity. All these options are compatible with the archived execution and analysis infrastructure and can be preregistered using the same evidence-recording conventions.

Threats to validity revisited (artifact-supported). Internal validity: by design the shared\_initial\_candidate semantics isolates detection but prevents verifier-mediated candidate resampling from influencing outcomes; this tradeoff explains why a degenerate result may occur when baseline error prevalence is low. Construct validity: emulator-observed faults are a conservative proxy for operational harm but do not capture traffic-dependent timing, vendor-specific hardware idiosyncrasies, or live-traffic emergent failures. External validity: experiments used a single declared model (gpt-5-mini), a synthetic task bank, and one deterministic verifier; extrapolation to other models, vendor dialects, or production hardware is limited. Conclusion validity: because baseline error prevalence was zero in this corpus the confirmatory design had no power to detect verifier benefit; the preregistered power target (detect absolute paired reduction 0.20) becomes moot under a zero-prevalence baseline.

Operational implications and measured tradeoffs. The confirmed ceiling outcome does not imply verifiers are unnecessary; rather, it shows that under certain tightly constrained synthesis and evaluation conditions an LLM alone produced valid device sequences on the synthetic corpus. In production settings with ambiguous intents, incomplete constraints, or vendor heterogeneity, verifiers remain a mechanistic guard. Measuring verifier costs (false positives that block safe changes) and benefits (true positives preventing faults) requires task banks with non-negligible baseline failure prevalence; the archived artifacts and preregistration provide a reproducible template for such follow-on evaluations.

Reproducibility notes (artifact pointers and deterministic checks). All execution and analysis artifacts are archived in the evidence bundle. Key artifacts referenced in this paper include: execution/raw\_results.jsonl, execution/task\_manifest.jsonl, execution/model\_configuration.json, analysis/paired\_contingency\_table.csv, analysis/condition\_summary.csv, analysis/analysis\_log.jsonl, deterministic\_reconciliation.json, and per-episode validator traces under execution/scoring/. The analysis bootstrap used seed = 1729 and 10,000 resamples; analysis/analysis\_log.jsonl records these values. The run-level timing, provider-call audit, and artifact counts are recorded in execution/ and analysis/ manifests to enable deterministic reconciliation of the paired accounting. The model configuration records temperature = null (execution/model\_configuration.json), which we report as provider-default sampling semantics (a residual sampling-parameter uncertainty archived in the evidence).

\section{Limitations}
\begin{itemize}
\item Confirmatory ceiling/null outcome: the baseline generator produced zero emulator faults on the preregistered task set, preventing direct measurement of verifier detection benefit.
\item Shared\_initial\_candidate semantics: the design measures verifier effect on the same generated candidate only and does not assess independent per-condition generation, multi-sample generation, or repairs unless explicitly permitted by the contract.
\item Single-model scope: experiments used one declared model (gpt-5-mini); results do not imply generality across models or versions.
\item Synthetic/emulated tasks: task corpus and emulator executions differ from production devices and live traffic; external validity to production deployments is limited.
\item Sampling-parameter uncertainty: execution/model\_configuration.json shows temperature = null (sampling parameters unset); null indicates provider-default runtime sampling semantics and is a residual uncertainty recorded in the archived configuration.
\item Residual contamination risk: deterministic provenance and exact-match checks were applied, but private training-data contamination of the hosted model cannot be fully ruled out and remains residual uncertainty.
\end{itemize}

\section{Conclusion}
We executed a preregistered paired evaluation (c1-gnr-vv-2026-0001) under shared\_initial\_candidate semantics to test whether a canonical deterministic verifier reduces the proportion of LLM-generated configurations that would execute with operational faults. For the chosen model (gpt-5-mini), verifier (deterministic\_netops\_validator\_v5), and synthetic/emulated task bank, baseline executions produced zero emulator faults and the verifier did not change outcomes (paired difference = 0.0; bootstrap 95\% CI [0.0,0.0]; exact McNemar p = 1.0). This artifact-grounded null/ceiling outcome is internally consistent with the preregistered design: when identical generated candidates produce no executed faults a verifier cannot reduce executed faults under the shared\_initial\_candidate contract. Measuring verifier utility under realistic error regimes requires preregistered contracts that permit independent or multiple generator samples, repair-enabled flows, adversarial task sets, and multi-model comparisons. The archived artifacts (responses, task manifests, validator traces, execution manifest, and analysis logs) provide a reproducible baseline and a template for those follow-on confirmatory designs and power calculations.

\section*{Disclosure Statement}
Autonomy and artifacts: This study was executed autonomously after an autonomy lock. Human role: a Research Director selected the topic family and approved the immutable initial master prompt prior to the autonomy lock; no human manually edited technical manuscript content after the lock. Models and tools: initial generation used gpt-5-mini (declared\_model\_version recorded in execution/model\_configuration.json) orchestrated via the hosted\_netops\_gvr\_v1 adapter; deterministic\_netops\_validator\_v5 (validator\_version = "5.0") served as the canonical verifier; an isolated emulator executed candidate configurations. Preregistration: study c1-gnr-vv-2026-0001 was preregistered and frozen before observing outcomes. Immutable master prompt reference: provenance/master\_prompt.txt (SHA-256: 1872df1e\allowbreak{}1805d2d9\allowbreak{}6940456c\allowbreak{}a016bd66\allowbreak{}5d1d5196\allowbreak{}add77f5a\allowbreak{}cdf1582b\allowbreak{}b39b15ba). Key archived artifacts include execution/raw\_results.jsonl, execution/task\_manifest.jsonl, execution/model\_configuration.json, analysis/paired\_contingency\_table.csv, analysis/condition\_summary.csv, deterministic\_reconciliation.json, and per-episode validator traces under execution/scoring/. The model configuration records temperature = null (provider-default sampling semantics archived in execution/model\_configuration.json). Provider-call audit: hosted\_model\_call\_event\_count = 72. No human manual manuscript editing or content insertion occurred after the autonomy lock.

\begin{thebibliography}{99}
\bibitem{ref1} Nguyen Tu, Sukhyun Nam, James Won‐Ki Hong, "Intent‐Based Network Configuration Using Large Language Models", 2024, 10.1002/nem.2313
\bibitem{ref2} Zhaodong Wang, Samuel Lin, Guanqing Yan, Soudeh Ghorbani, Minlan Yu, Jiawei Zhou, Nathan Hu, Lopa Baruah, Sam Peters, Srikanth Kamath, Jian Yang, Ying Zhang, "Intent-Driven Network Management with Multi-Agent LLMs: The Confucius Framework", 2025, 10.1145/3718958.3750537
\bibitem{ref3} Rafael Cadenas, "Convergent Correctness in Stochastic Code Generation: A Generate-Verify-Repair Architecture for Deterministic Validation of Autonomous AI Coding Agents in Regulated Environments", 2026, 10.2139/ssrn.6754899
\bibitem{ref4} http://citeseerx.ist.psu.edu/viewdoc/summary?doi=10.1.1.300.8659
\end{thebibliography}

\end{document}
